\begin{abstract}
Vision Transformers (ViTs) incur quadratic cost with input resolution, making high-resolution training expensive. Resolution extrapolation---training at low resolution and inferring at higher resolution---offers a promising alternative, yet standard positional embeddings fail to generalize beyond the training grid. While relative positional encodings have enabled context extrapolation in large language models, their effectiveness in the 2D spatial domain remains underexplored.

We present a systematic evaluation of modern positional encoding methods adapted for Vision Transformers, comparing learned absolute embeddings, RoPE, ALiBi, YaRN, and FIRE across classification, detection, and segmentation tasks. Evaluating zero-shot extrapolation from $224 \times 224$ up to $1024 \times 1024$ ($4.6\times$ scale), we find that attention bias-based methods substantially outperform rotation-based alternatives: FIRE retains $63.1\%$ Top-1 accuracy at $1024$ px versus RoPE's $18.9\%$, a $44$ percentage point gap. These trends persist in dense prediction, where ALiBi and FIRE maintain stable performance while RoPE degrades sharply. Our attention analysis reveals that bias-based methods preserve focused, semantically meaningful attention patterns under extrapolation, whereas rotation-based methods exhibit attention collapse. These findings provide practical guidance for building resolution-flexible Vision Transformers under limited compute budgets. Reproducible code and weights will be released.
\end{abstract}